\subsection{Internal and External Links}

The \texttt{hyperref} package allows clickable links inside the document,
including references, the table of contents, URLs, and links to other
sections.

\begin{lstlisting}[style=latex]
\usepackage{hyperref}

\hypersetup{
    colorlinks=true,
    linkcolor=blue,      % Internal links
    filecolor=magenta,
    urlcolor=cyan,
    linktoc=all          % Title and page number are clickable
}
\end{lstlisting}

To link to another part of the document, first create a label:

\begin{lstlisting}[style=latex]
\section{Control Systems}
\label{sec:control}
\end{lstlisting}

Then reference it using:

\begin{lstlisting}[style=latex]
See Section~\ref{sec:control}.
\end{lstlisting}

A clickable text link can also be created with:

\begin{lstlisting}[style=latex]
\hyperref[sec:control]{Go to Control Systems}
\end{lstlisting}

For external URLs:

\begin{lstlisting}[style=latex]
\url{https://example.com}

\href{https://example.com}{Example website}
\end{lstlisting}

If visible link colours or borders are not desired, use:

\begin{lstlisting}[style=latex]
\usepackage[hidelinks]{hyperref}
\end{lstlisting}

With \texttt{hidelinks}, the links remain clickable, but no coloured
text or link borders are shown.